\documentclass{article}
\usepackage{graphicx} % Required for inserting images
\usepackage{setspace}
\usepackage{amsmath}
\usepackage{float}


\title{\textbf{Transformasi Geometri}}


\begin{document}
\onehalfspacing
\maketitle

Transformasi yang dimaksud dalam materi ini adalah perpindahan atau perubahan ukuran suatu objek yang disebabkan oleh pergeseran, pencerminan, perputaran, dan penskalaan.

\section{Translasi}
Translasi adalah pergeseran suatu objek tanpa mengubah bentuk dan ukurannya. Dalam materi ini objek yang akan di translasikan berupa titik, garis, ataupun bidang dalam koordinat kartesius. Misalnya saja titik $A(2,3)$, jika titik tersebut digeser oleh $\begin{bmatrix}
    6\\-7
\end{bmatrix}$ maka akan menghasilkan
\[
\begin{aligned}
A'=
    \begin{bmatrix}
        2\\-3
    \end{bmatrix}+\begin{bmatrix}
        6\\-7
    \end{bmatrix}=\begin{bmatrix}
        8\\-10
    \end{bmatrix}
\end{aligned}
\]
dimana $A'$ disebut bayangan dari $A$.\\
Diketahui garis $g:2y=5x+7$. Jika garis $g$ ditranslasikan oleh $\begin{bmatrix}
    3\\-4
\end{bmatrix}$ maka garis baru yang akan dihasilkan adalah\\
\[\begin{aligned}
    x&=x'-3\\
    y&=y'-(-4)=y'+4\\
    2y&=5x+7\\
    2(y'+4)&=5(x'-3)+7\\
    2y'+8&=5x'-15+7\\
    2y'&=5x'-15+7-8\\
    2y'&=5x'-16
\end{aligned}\]
Jadi garis baru yang dihasikan adalah $2y=5x-16$ atau bisa ditulis $\displaystyle y=\frac{5}{2}x-8$.\\
\textbf{Contoh Lain}. Diketahui titik $A(8,1),B(12,1),C(10,5)$ membentuk sebuah segitiga. Jika segitiga ABC ditranslasikan oleh $\begin{bmatrix}
    0\\-3
\end{bmatrix}$, maka luas bagian segitiga yang ada di kuadran II adalah $\dots$
\[
\begin{aligned}
    A'&=\begin{bmatrix}
        8\\1
    \end{bmatrix}+\begin{bmatrix}
        0\\-3
    \end{bmatrix}=\begin{bmatrix}
        8\\-2
    \end{bmatrix}\\
    B'&=\begin{bmatrix}
        12\\1
    \end{bmatrix}+\begin{bmatrix}
        0\\-3
    \end{bmatrix}=\begin{bmatrix}
        12\\-2
    \end{bmatrix}\\
    C'&=\begin{bmatrix}
        10\\5
    \end{bmatrix}+\begin{bmatrix}
        0\\-3
    \end{bmatrix}=\begin{bmatrix}
        10\\2
    \end{bmatrix}\\
\end{aligned}
\]
Akan coba digambar segitiga yang telah di translasikan tersebut
\begin{figure}[H]
    \centering
    \includegraphics[width=1\linewidth]{gambar/bayangansegitiga.png}
\end{figure}
Dari gambar diatas, dapat dihitung luas bagian segitiga yang ada di kuadran II adalah
\[
\begin{aligned}
    L&=\frac{1}{2}\times 4 \times4-\frac{1}{2}\times2\times2\\
    L&=8-2=6
\end{aligned}
\]
Jadi luas bagian segitiga yang ada di kuadran II adalah 6 satuan luas.
\section{Refleksi}
\section{Rotasi}
\section{Dilatasi}
\end{document}
